\documentclass[11pt]{article}

\usepackage[a4paper,margin=1in]{geometry}
\usepackage{amsmath,amssymb,amsthm,mathtools}
\usepackage[T1]{fontenc}
\usepackage{lmodern}
\usepackage{microtype}
\usepackage[hidelinks]{hyperref}

\newtheorem{theorem}{Theorem}
\newtheorem{proposition}[theorem]{Proposition}
\newtheorem{corollary}[theorem]{Corollary}
\theoremstyle{remark}
\newtheorem{remark}[theorem]{Remark}

\newcommand{\A}{\mathbb A}
\newcommand{\Z}{\mathbb Z}

\title{A Four-Parameter Polynomial Parametrization for
\texorpdfstring{\(xyz+ts+1=0\)}{xyz+ts+1=0}}
\author{}
\date{}

\begin{document}

\maketitle
\vspace{-2.5em}

\begin{abstract}
We give an explicit four-parameter polynomial family of integer solutions to
\(xyz+ts+1=0\).  The construction comes from a two-factor identity for
\(1+xyz\).  Over any field, the resulting map is an isomorphism on the open
set where \(xy\neq0\); thus it is a birational parametrization with the
expected number of parameters.  We also characterize exactly which integral
points are obtained from integral parameters.
\end{abstract}

\section{The polynomial family}

\begin{theorem}\label{thm:main}
Fix \(\varepsilon\in\{1,-1\}\).  For arbitrary parameters
\(a,b,m,n\in\Z\), define
\begin{equation}\label{eq:param-epsilon}
\begin{aligned}
 x&=a, &\qquad y&=b,\\
 z&=\varepsilon(m+n)+abmn, &
 t&=\varepsilon+abm,\\
 &&s&=-\varepsilon-abn.
\end{aligned}
\end{equation}
Then
\[
xyz+ts+1=0.
\]
Equivalently, the five polynomials in \eqref{eq:param-epsilon} satisfy the
identity \(xyz+ts+1\equiv0\) in \(\Z[a,b,m,n]\).
\end{theorem}

\begin{proof}
Put \(q=ab\).  Since \(\varepsilon^2=1\),
\begin{equation}\label{eq:keyidentity}
(\varepsilon+qm)(\varepsilon+qn)
 =1+\varepsilon q(m+n)+q^2mn
 =1+q\bigl(\varepsilon(m+n)+qmn\bigr).
\end{equation}
Under \eqref{eq:param-epsilon}, the last expression is \(1+xyz\), whereas
\[
ts=-(\varepsilon+qm)(\varepsilon+qn).
\]
Thus \(ts=-(1+xyz)\), which is equivalent to the required equation.
\end{proof}

Taking \(\varepsilon=1\) gives the particularly simple parametrization
\begin{equation}\label{eq:main-param}
\boxed{
(x,y,z,t,s)
 =\bigl(a,b,m+n+abmn,\,1+abm,\,-1-abn\bigr),
\qquad a,b,m,n\in\Z.
}
\end{equation}
For instance, setting \(a=b=1\), \(m=0\), and \(n=r\) yields the infinite
subfamily
\[
(x,y,z,t,s)=(1,1,r,1,-1-r),\qquad r\in\Z.
\]

\section{How the construction was found}

The equation may be rewritten as
\begin{equation}\label{eq:factor-target}
1+xyz=-ts.
\end{equation}
Write \(q=xy\).  A direct way to enforce \eqref{eq:factor-target} is to make
\(1+qz\) split into two polynomial factors.  Introduce two free parameters
\(m,n\) and seek a factorization of the form
\[
1+qz=(\varepsilon+qm)(\varepsilon+qn),
\qquad \varepsilon^2=1.
\]
Expanding the right-hand side gives
\[
1+\varepsilon q(m+n)+q^2mn,
\]
so the factorization holds precisely when
\[
z=\varepsilon(m+n)+qmn.
\]
One may then take one factor to be \(t\) and the negative of the other to be
\(s\):
\[
t=\varepsilon+qm,
\qquad
s=-\varepsilon-qn.
\]
Finally, writing \(q=ab\) leaves \(x=a\) and \(y=b\) as independent
parameters.  This produces exactly \eqref{eq:param-epsilon}.

For \(\varepsilon=1\), the central identity can also be expressed through the
binary operation
\[
m\mathbin{\oplus_q}n:=m+n+qmn.
\]
It is designed so that
\[
1+q(m\mathbin{\oplus_q}n)=(1+qm)(1+qn).
\]
Thus the nonlinear-looking term \(m+n+qmn\) is not guessed: it is forced by
multiplying two expressions congruent to \(1\) modulo \(q\).

\section{Generic completeness}

Let
\[
V:\ xyz+ts+1=0
\]
be the affine hypersurface in \(\A^5\), and let \(\Phi_\varepsilon:\A^4\to V\)
be the polynomial map defined by \eqref{eq:param-epsilon}.  For every field
\(K\), the polynomial \(xyz+ts+1\) is irreducible in
\(K[x,y,z,t,s]\): viewed as a polynomial in \(z\), it is primitive over
\(K[x,y,t,s]\) and linear over that ring's fraction field.  Hence the
nonempty open subset \(xy\neq0\) is dense in \(V\).

\begin{proposition}\label{prop:birational}
Over any field \(K\), the restriction of \(\Phi_\varepsilon\) to \(ab\neq0\)
is an isomorphism onto the open subset of \(V\) on which \(xy\neq0\).  Its
inverse is
\begin{equation}\label{eq:inverse}
 a=x,
 \qquad b=y,
 \qquad m=\frac{t-\varepsilon}{xy},
 \qquad n=-\frac{s+\varepsilon}{xy}.
\end{equation}
In particular, \(\Phi_\varepsilon\) is a birational parametrization of \(V\).
\end{proposition}

\begin{proof}
Let \((x,y,z,t,s)\in V(K)\) with \(q:=xy\neq0\), and define
\(a,b,m,n\) by \eqref{eq:inverse}.  The formulas for \(x,y,t,s\) in
\eqref{eq:param-epsilon} are then immediate.  Moreover,
\begin{align*}
q\bigl(\varepsilon(m+n)+qmn\bigr)
&=\varepsilon(t-\varepsilon)-\varepsilon(s+\varepsilon)
  -(t-\varepsilon)(s+\varepsilon)\\
&=-ts-1.
\end{align*}
Since the point lies on \(V\), one has \(-ts-1=xyz=qz\).  Division by
\(q\) therefore gives
\[
z=\varepsilon(m+n)+qmn.
\]
Hence \eqref{eq:inverse} is indeed the inverse map.  Both maps are regular on
the indicated principal open subsets.
\end{proof}

The proposition explains why four independent parameters arise naturally:
one equation in five affine variables has four-dimensional generic freedom,
and the construction realizes that freedom without redundancy on
\(xy\neq0\).

\section{Integral image and the boundary
\texorpdfstring{\(xy=0\)}{xy=0}}

The preceding proposition is a completeness statement over a field.  For
integer parameters, divisibility enters because the inverse formulas in
\eqref{eq:inverse} need not be integral.

\begin{proposition}\label{prop:integer-image}
Let \((x,y,z,t,s)\in\Z^5\) satisfy \(xyz+ts+1=0\), and suppose \(xy\neq0\).
For fixed \(\varepsilon\in\{1,-1\}\), this point belongs to
\(\Phi_\varepsilon(\Z^4)\) if and only if
\begin{equation}\label{eq:congruence}
 t\equiv\varepsilon\pmod{xy}.
\end{equation}
Equivalently, one may require \(s\equiv-\varepsilon\pmod{xy}\).
\end{proposition}

\begin{proof}
If the point is obtained from integral parameters, then
\(t-\varepsilon=xy\,m\) and \(s+\varepsilon=-xy\,n\), so both congruences
hold.

Conversely, assume \eqref{eq:congruence}.  Reducing the original equation
modulo \(xy\) gives \(ts\equiv-1\pmod{xy}\).  Since
\(t\equiv\varepsilon\) and \(\varepsilon^2=1\), it follows that
\(s\equiv-\varepsilon\pmod{xy}\).  Hence the values of \(m,n\) in
\eqref{eq:inverse} are integers, and Proposition~\ref{prop:birational} shows
that these parameters recover the given point.
\end{proof}

If \(xy=0\), the equation reduces to \(ts=-1\).  Hence an integral solution
has either \((t,s)=(1,-1)\) or \((t,s)=(-1,1)\), while \(z\) is arbitrary.
The family \(\Phi_1\) covers the first case and \(\Phi_{-1}\) covers the
second.  Thus the two sign choices cover the entire degenerate integral locus
\(xy=0\).

\begin{remark}
The formulas above give a polynomial family of integral solutions and a
birational parametrization of the full hypersurface.  They are not asserted
to send integral parameters onto every integral point with \(xy\neq0\);
Proposition~\ref{prop:integer-image} gives the exact congruence criterion for
such integral lifting.
\end{remark}

\end{document}
